\documentclass[conference]{IEEEtran}
\usepackage{cite}
\usepackage{graphicx}
\usepackage{amsmath,amssymb}
\usepackage{algorithmic}
\usepackage{array}
\usepackage{booktabs}

% For better table spacing
\setlength{\parindent}{0pt}
\setlength{\parskip}{0.5ex}

\begin{document}

% Title
\title{Reinforcement Learning for Adaptive Neighborhood Selection in Vehicle Routing with Synchronization Constraints}

% Author block 
\author{
    \IEEEauthorblockN{Nada Farah Djedjig}
    \IEEEauthorblockA{\textit{Department of Computer Science} \\
        \textit{University Name}\\
        City, Country \\
        nada.djedjig@example.com}
    \and
    \IEEEauthorblockN{Supervisor Name}
    \IEEEauthorblockA{\textit{Department} \\
        \textit{University Name}\\
        City, Country \\
        supervisor@example.com}
    \and
    \IEEEauthorblockN{Supervisor Name}
    \IEEEauthorblockA{\textit{Department} \\
        \textit{University Name}\\
        City, Country \\
        supervisor@example.com}
}

\maketitle

\begin{abstract}
The Vehicle Routing Problem with Time Windows and Synchronization constraints (VRPTW‑S) is NP‑hard and central to home healthcare logistics. While metaheuristics such as GVNS achieve state‑of‑the‑art performance, they rely on a fixed order of neighborhood operators, which cannot adapt to the changing search context. This paper presents the first systematic comparison of two reinforcement learning agents – tabular Q‑learning and Deep Q‑Network (DQN) – for dynamic operator selection in VRPTW‑S. We formulate the selection task as a Markov Decision Process and integrate both agents into a GVNS framework. Experiments on Solomon‑based instances (50–400 customers) show that Q‑learning achieves up to 72\% improvement on small instances (50 customers) but fails to scale, while DQN remains stable and reaches 83\% improvement on large instances (400 customers). Our results demonstrate that DQN’s ability to generalize across continuous states makes it the preferred choice for large‑scale home healthcare routing.
\end{abstract}

\begin{IEEEkeywords}
Vehicle Routing Problem, Time Windows, Synchronization, Reinforcement Learning, Q-learning, Deep Q-Network, Home Healthcare.
\end{IEEEkeywords}

\section{INTRODUCTION}
\label{sec:intro}

\begin{itemize}
    \item \textbf{Context:} Aging population and chronic diseases drive home healthcare (HHC) demand → VRPTW‑S as key logistics problem. (Source: Chap. 1, §1.1)
    \item \textbf{Challenge:} Synchronization + time windows create combinatorial explosion; exact methods limited to $\approx$20 customers. (Chap. 1, §1.3)
    \item \textbf{Gap:} State‑of‑the‑art metaheuristics (GVNS, ALNS) use fixed operator order; cannot adapt to search dynamics. (Chap. 4, §4.5)
    \item \textbf{Contribution:}
        \begin{itemize}
            \item First comparison of tabular Q‑learning vs DQN for operator selection in VRPTW‑S.
            \item MDP formulation with compact state features (instance, solution, progress, operator success rates).
            \item Experimental evaluation on 50–400 customer instances.
            \item Key insight: QL excels on small scales; DQN scales and stabilizes.
        \end{itemize}
    \item \textbf{Paper structure:} Section II – background; Section III – MDP formulation; Section IV – QL and DQN agents; Section V – experiments; Section VI – conclusion.
\end{itemize}

\section{BACKGROUND}
\label{sec:background}

\subsection{VRPTW‑S}
\label{subsec:vrptws}
\begin{itemize}
    \item Definition: graph $G=(V,E)$ with time windows $[e_i,l_i]$, service times $s_i$, synchronization constraints (strict, overlap, precedence). (Chap. 1, §1.2)
    \item Complexity: NP‑hard; exact methods only for tiny instances. (Chap. 1, §1.4)
    \item Baseline metaheuristic: GVNS (shaking + VND + penalty self‑adjustment) – see Sec. IV. (Chap. 4)
\end{itemize}

\subsection{Reinforcement Learning Fundamentals}
\label{subsec:rl}
\begin{itemize}
    \item Markov Decision Process (MDP): $(S,A,P,R,\gamma)$; agent learns policy $\pi(a|s)$ to maximize cumulative reward. (Chap. 3, §3.2)
    \item Value‑based methods: Q‑learning (tabular), DQN (neural approximation). (Chap. 3, §3.4)
    \item Exploration–exploitation trade‑off: UCB, $\epsilon$-greedy. (Chap. 5, §5.3.2)
\end{itemize}

\section{MDP FORMULATION FOR OPERATOR SELECTION}
\label{sec:mdp}

\begin{itemize}
    \item \textbf{State space $S$:} 24‑dimensional feature vector (normalised):
        \begin{itemize}
            \item Instance level: customer count, time‑window width, synchronisation ratio.
            \item Solution quality: penalised cost, travel time, penalties.
            \item Search progress: stagnation level, elapsed time ratio.
            \item Neighbourhood success rates (8 values, rolling window).
        \end{itemize}
        (Chap. 5, §5.2.1)
    \item \textbf{Action space $A$:} 8 neighbourhood operators:
        \begin{itemize}
            \item Intra‑route insertion, inter‑route insertion, intra‑route swap, inter‑route swap,
            \item $2$-opt, $2$-opt$^*$, cross exchange, $k$-relocate.
        \end{itemize}
        (Chap. 4, §4.3)
    \item \textbf{Reward $R$:} Normalised improvement $\Delta F_{\text{norm}}$, shaped with bonus for feasibility and penalty for time. (Chap. 5, §5.2.3)
    \item \textbf{Transition $P$:} Deterministic – applying an operator yields a new solution and next state.
\end{itemize}

\section{METHODOLOGY: Q-LEARNING AND DQN AGENTS}
\label{sec:agents}

\subsection{Tabular Q-Learning}
\label{subsec:ql}
\begin{itemize}
    \item State discretisation → 144 abstract states (feasibility flags, stagnation, effectiveness, instance size). (Chap. 5, §5.3.1)
    \item Q‑table $Q(s,a)$, updated by Bellman equation:
        \[
        Q(s,a) \leftarrow Q(s,a) + \alpha \left[ r + \gamma \max_{a'} Q(s',a') - Q(s,a) \right]
        \]
    \item Action selection: UCB (Upper Confidence Bound) for balanced exploration. (Chap. 5, §5.3.2)
    \item Hyperparameters: $\alpha=0.1$, $\gamma=0.95$, $c_{\text{UCB}}=1.0$.
\end{itemize}

\subsection{Deep Q-Network (DQN)}
\label{subsec:dqn}
\begin{itemize}
    \item Neural approximator: MLP with input layer (24 features), two hidden layers (128 neurons each, ReLU), output layer (8 Q‑values). (Chap. 5, §5.4.1)
    \item Online Q‑network $Q_\theta$, target network $Q_{\theta^-}$ (updated every 100 steps).
    \item Experience replay buffer (size 50,000); mini‑batch size 32.
    \item $\epsilon$-greedy exploration ($\epsilon$ decays from 1.0 to 0.05 over 500 steps).
    \item Learning rate $3\times10^{-4}$, Huber loss.
\end{itemize}

\subsection{Integration into GVNS}
\begin{itemize}
    \item At each neighbourhood evaluation, the RL agent selects one of 8 operators.
    \item The selected operator is applied; reward and next state are computed.
    \item The agent updates its Q‑table / network parameters.
    \item (Fig. 19 in report shows the architecture.)
\end{itemize}

\section{EXPERIMENTAL RESULTS}
\label{sec:results}

\subsection{Experimental Setup}
\begin{itemize}
    \item \textbf{Instances:} Solomon (C, R, RC) with 50, 100, 200, 400 customers; synchronization added (5/10/20/40 pairs). (Chap. 4, §4.6)
    \item \textbf{Baseline:} Classical GVNS with fixed operator order.
    \item \textbf{Metrics:} % gap = $( \text{obj}_{\text{RL}} - \text{obj}_{\text{GVNS}} ) / \text{obj}_{\text{GVNS}} \times 100$; 5 runs per instance.
\end{itemize}

\subsection{Q-Learning Results}
\begin{itemize}
    \item \textbf{50 customers:} outperforms GVNS in $\approx$72\% of instances; gaps $-0.1\%$ to $-2\%$. (Chap. 5, §5.3.2)
    \item \textbf{100 customers:} $\approx$55\% improvement; more variable (up to $+5.5\%$). (Chap. 5, §5.3.2)
    \item \textbf{200–400 customers:} fails to scale; worst gap $+24.8\%$ on 400 customers. (Chap. 5, §5.3.2)
\end{itemize}

\subsection{DQN Results}
\begin{itemize}
    \item \textbf{50 customers:} $\approx$45\% improvement (less consistent than QL). (Chap. 5, §5.4.2)
    \item \textbf{100 customers:} $\approx$67\% improvement; gaps up to $-2.5\%$, stable. (Chap. 5, §5.4.2)
    \item \textbf{200–400 customers:} $\approx$83\% improvement; gaps $-0.5\%$ to $-1.8\%$, reliable. (Chap. 5, §5.4.2)
\end{itemize}

\subsection{Comparison and Discussion}
\begin{itemize}
    \item \textbf{Key insight:} QL is simple and effective for small problems ($n\le100$) but collapses due to state space explosion. DQN generalises via neural approximation and handles large instances robustly.
    \item \textbf{Runtime:} DQN often finds good solutions faster than GVNS; QL overhead is moderate.
    \item \textbf{Takeaway for HHC:} For real‑world large‐scale systems, DQN‐based adaptive operator selection is the preferred choice.
\end{itemize}

\section{CONCLUSION AND FUTURE WORK}
\label{sec:conclusion}

\begin{itemize}
    \item We presented the first systematic comparison of Q‑learning and DQN for operator selection in VRPTW‑S.
    \item QL excels on small instances; DQN scales to 400+ customers with consistent improvements.
    \item Future work: integrate graph neural networks for richer state representation; extend to multi‑objective optimisation (cost, equity, continuity of care).
    \item \textbf{Link to CIE53 theme:} Adaptive routing reduces travel distances and fuel consumption, contributing to sustainable healthcare logistics.
\end{itemize}

\section*{Acknowledgment}
The authors thank ... (optional)

\begin{thebibliography}{9}
\bibitem{bre08} D. Bredström and M. Rönnqvist, ``Combined vehicle routing and scheduling with temporal precedence and synchronization constraints,'' \textit{EJOR}, vol. 191, no. 1, pp. 19–31, 2008.
\bibitem{mas25} M. Masmoudi, R. Borchani, and B. Jarboui, ``Generalized variable neighborhood search algorithm for vehicle routing problem with time windows and synchronization,'' \textit{Computers \& Operations Research}, vol. 183, 107193, 2025.
\bibitem{liu19} R. Liu, Y. Tao, and X. Xie, ``An adaptive large neighborhood search heuristic for the vehicle routing problem with time windows and synchronized visits,'' \textit{Computers \& Operations Research}, vol. 101, pp. 250–262, 2019.
\bibitem{sut18} R. S. Sutton and A. G. Barto, \textit{Reinforcement Learning: An Introduction}, 2nd ed. MIT Press, 2018.
\bibitem{mnih15} V. Mnih et al., ``Human-level control through deep reinforcement learning,'' \textit{Nature}, vol. 518, no. 7540, pp. 529–533, 2015.
\bibitem{bel17} I. Bello, H. Pham, Q. V. Le, M. Norouzi, and S. Bengio, ``Neural combinatorial optimization with reinforcement learning,'' in \textit{ICLR}, 2017.
\bibitem{kool19} W. Kool, H. van Hoof, and M. Welling, ``Attention, learn to solve routing problems!'' in \textit{ICLR}, 2019.
\bibitem{solo87} M. M. Solomon, ``Algorithms for the vehicle routing and scheduling problems with time window constraints,'' \textit{Operations Research}, vol. 35, no. 2, pp. 254–265, 1987.
\end{thebibliography}

\end{document}